\documentclass[10pt,a4paper]{amsart}
\usepackage[margin=19mm]{geometry}
\usepackage{amsmath,amssymb,mathtools}
\usepackage[scr=rsfs,cal=euler]{mathalfa}
\usepackage{xfrac}
\usepackage[unicode,hidelinks,bookmarks=false]{hyperref}
\newcommand{\ie}{i.e.\ }
\newcommand{\cf}{cf.\ }
\newcommand{\resp}{respectively\ }
\newcommand{\Subsection}{\subsection}
\providecommand{\nobreakdash}{\nobreak\hbox{-}}
\setlength{\emergencystretch}{2em}
\multlinegap0pt
\usepackage{amssymb}
\usepackage{mathtools}
\newtheorem{theorem}{Theorem}
\newcommand{\ac}{_\mathrm{ac}}
\newcommand{\bbR}{\mathbb{R}}
\newcommand{\dc}{_\mathrm{dc}}
\newcommand{\dif}{\mathrm{d}}
\newcommand{\mean}[1]{\langle #1 \rangle}
\DeclareMathOperator{\median}{med}
\newcommand{\norm}[1]{\Vert #1 \Vert}
\newcommand{\normx}[1]{\left\Vert #1 \right\Vert}
\newcommand{\pb}[1]{\big( #1 \big)}
\newcommand{\trans}{^\mathsf{T}}
\title{\LARGE \bf Period-Aware Asymptotic Gain\\with Application to a Periodically Forced Synchronization Circuit}
\author{Anton Ponomarev, Lutz Gr\"oll, Veit Hagenmeyer}
\date{Source: arXiv:2603.24600v1 (CC BY 4.0)}
\begin{document}
\maketitle

\noindent\emph{Source and licence.} \LARGE \bf Period-Aware Asymptotic Gain\\with Application to a Periodically Forced Synchronization Circuit. Version arXiv:2603.24600v1 (CC BY 4.0), distributed by arXiv under the Creative Commons Attribution 4.0 International licence, http://creativecommons.org/licenses/by/4.0/, as recorded in the arXiv licence metadata for this exact version and shown on the abstract page https://arxiv.org/abs/2603.24600v1. Full licence notice: \texttt{licenses/arxiv-2603.24600-CC-BY-4.0.txt}.\par\smallskip

\noindent\textit{Editorial scope.} Counted unit: the article's theorem linear (source0.tex lines 354--373). The article's complete original proof of it is delivered with it (source0.tex lines 375--415, one \texttt{proof} environment of 41 lines). No mathematics is written, repaired or re-derived: the statement and the proof are the article's own lines, in the article's own notation, and no retained line is substituted or re-flowed.  Retained uncounted as the source-local context the proof needs: lines 199--202.  Nothing was omitted from any retained span, so the package carries no omission record.  No retained line contains a citation, so the wrapper carries no bibliography and no bibliography entry.  No retained line contains a figure environment, an image-inclusion command or drawing code, so no image is delivered and the package declares no asset dependency.  Editorial formatting only, in addition: an AMS \texttt{amsart} preamble that restates the article's own declarations and macros used by the retained lines, namely the environments \texttt{theorem} on separate counters, together with the standard packages those lines need.  The retained environments print in this document as Theorem 1 in order of appearance, whereas the article prints the same items with the numbering of its own full document.  The complete native source of the article as distributed in the arXiv e-print for this exact version is bundled unchanged as \texttt{sources/197156/source0.tex}.
    \begin{equation}
        \label{eq: median definition implicit}
        \frac1T \int_0^T \frac{f(t) - \median f}{\norm{f(t) - \median f}} \,\dif t = 0
    \end{equation}

\begin{theorem}
    \label{th: linear}
    \begin{subequations}
        The exact PAG of a stable linear system with transfer matrix $G$ and $T$-periodic-impulse response $H_T$ is
        \begin{equation}
            \gamma_T\pb{\rho_T(u)} = \Gamma_T \rho_T(u)
        \end{equation}
        where
        \begin{align}
            \Gamma_T &= \begin{bmatrix}
                \gamma\dc & 0 \\
                0 & \gamma\ac(T)
            \end{bmatrix}, \\
            \label{eq: dc gain linear}
            \gamma\dc &= \norm{G(0)}, \\
            \label{eq: ac gain linear exact}
            \gamma\ac(T) &= T \sup_{\substack{v\in\bbR^p \\ \norm{v} = 1}} D_{\median} (v \cdot H_T).
        \end{align}
    \end{subequations}
\end{theorem}

\begin{proof}
    Since $A$ is Hurwitz, the system's response to a $T$-periodic input $u(t) = u\dc + u\ac(t)$ globally converges to the $T$-periodic output $\hat y(t) = \hat y\dc + \hat y\ac(t)$ with
    \begin{subequations}
        \begin{align}
            \label{eq: dc linear transform}
            \hat y\dc &= G(0) u\dc,\\
            \label{eq: ac linear transform}
            \hat y\ac(t) &= \int_0^\infty H(\tau) u\ac(t-\tau) \,\dif\tau.
        \end{align}
    \end{subequations}
    From~\eqref{eq: dc linear transform} follows~\eqref{eq: dc gain linear}.
    
    Since $u\ac$ is $T$-periodic, \eqref{eq: ac linear transform} can be rearranged into the \emph{periodic convolution}
    \begin{equation}
        \label{eq: ac linear transform periodic}
        \hat y\ac(t) \equiv \int_0^T H_T(\tau) u\ac(t-\tau) \,\dif\tau.
    \end{equation}
    Thus,
    \begin{equation}
        \label{eq: output norm as supremum}
        \norm{\hat y\ac(t)} = \sup_{\substack{v\in\bbR^p \\ \norm{v} = 1}}
        \normx{\int_0^T v \cdot H_T(\tau) u\ac(t-\tau) \,\dif\tau}.
    \end{equation}
    If $\mean{u\ac}_T = 0$ then for every constant $\mu \in \bbR^p$
    \begin{equation}
        \norm{\hat y\ac(t)} \equiv \sup_{\substack{v\in\bbR^p \\ \norm{v} = 1}}
        \normx{\int_0^T \pb{v \cdot H_T(\tau) - \mu\trans} u\ac(t-\tau) \,\dif\tau}.
    \end{equation}
    The exact maximum of
    \begin{equation}
        \label{eq: ac integral to maximize}
        \normx{\int_0^T \pb{v \cdot H_T(\tau) - \mu\trans} u\ac(t-\tau) \,\dif\tau}
    \end{equation}
    with respect to $u\ac(\cdot)$ bounded by $\norm{u\ac}_\infty$ is delivered by
    \begin{equation}
        \label{eq: worst case input linear}
        u\ac(t-\tau) = \frac{v \cdot H_T(\tau) - \mu\trans}
        {\norm{v \cdot H_T(\tau) - \mu\trans}} \norm{u\ac}_\infty.
    \end{equation}
    In order to satisfy the requirement $\mean{u\ac}_T = 0$, it is necessary to take $\mu\trans = \median(v \cdot H_T)$ according to~\eqref{eq: median definition implicit}. With this $\mu$, the value of~\eqref{eq: ac integral to maximize} using~\eqref{eq: worst case input linear} is $T D_{\median}(v \cdot H_T) \norm{u\ac}_\infty$ which leads to~\eqref{eq: ac gain linear exact}.
\end{proof}

\end{document}
